\section{Conexión con DBeaver}

\subsection{Instalación del cliente}

Se descargó DBeaver Community 26.2.0 desde el sitio oficial. El archivo RPM se
verificó antes de instalarse y, posteriormente, se confirmó el paquete
registrado en Fedora.

\begin{lstlisting}[style=terminal]
$ curl -fL -o /tmp/dbeaver-ce-26.2.0-linux-x86_64.rpm \
    https://dbeaver.io/files/dbeaver-ce-latest-linux-x86_64.rpm
$ rpm -K /tmp/dbeaver-ce-26.2.0-linux-x86_64.rpm
$ sudo dnf install /tmp/dbeaver-ce-26.2.0-linux-x86_64.rpm
$ rpm -q dbeaver-ce
dbeaver-ce-26.2.0-stable.x86_64
\end{lstlisting}

Al abrir DBeaver por primera vez se descargó el controlador PostgreSQL JDBC
42.7.13. Después se creó una conexión con los parámetros de la
Tabla~\ref{tab:parametros-dbeaver}.

\begin{table}[H]
  \centering
  \rowcolors{2}{unamblue!6}{white}
  \begin{tabular}{ll}
    \toprule
    \textbf{Parámetro} & \textbf{Valor} \\
    \midrule
    Nombre de la conexión & PostgreSQL - Práctica 01 \\
    Servidor & \texttt{localhost} \\
    Puerto & \texttt{5432} \\
    Base de datos & \texttt{postgres} \\
    Usuario & \texttt{postgres} \\
    Contraseña & Credencial temporal oculta \\
    \bottomrule
  \end{tabular}
  \caption{Parámetros utilizados para la conexión local.}
  \label{tab:parametros-dbeaver}
\end{table}

La opción de conexión se ejecutó con la contraseña suministrada desde un
archivo de variables temporal. PostgreSQL confirmó tres contextos de DBeaver:
la conexión principal, la consulta de metadatos y la consola SQL.

\evidencia{05_evidencia.png}{Versión de DBeaver y sesiones del cliente observadas por PostgreSQL.}

La interfaz mostró la conexión activa en el navegador y abrió una consola
asociada al esquema \texttt{public} de la base \texttt{postgres}. Esta evidencia
confirma que el controlador JDBC, la autenticación, el puerto y el servidor
funcionaron de manera conjunta.

\evidencia[0.86\textwidth]{06_dbeaver_conectado.png}{DBeaver Community conectado al servidor PostgreSQL de la práctica.}
